\chapter{The Product}
\label{ch:product}

\section{Introduction}

This chapter presents the finished platform as a user experiences it. It walks
through the web application feature by feature, using screenshots of the running
system: connecting and processing libraries, asking grounded questions,
collaborating as a team, generating visuals and documents with the agent, saving
the resulting artefacts, and exploring the knowledge graph. Together these views
show that the design and implementation of the previous chapters come together
into a single, usable product.

\section{Libraries and Processing}

The starting point is the libraries view, shown in Figure~\ref{fig:ui-libraries}.
A user connects a source and creates a library, and the platform processes its
documents through the pipeline described in Chapter~5. Each library shows its
processing state and a progress bar, so the user can see at a glance when a
collection is ready to be used. The figure shows several processed libraries,
including collections of research papers, spreadsheets and source code, each
marked as complete.

\begin{figure}[H]
  \centering
  \includegraphics[width=\textwidth]{Preprocess_Libraries.png}
  \caption{The libraries view. A user creates a library from a connected source,
  and the platform processes its documents in the background. Each card shows the
  library's status and progress, here several completed collections of papers,
  spreadsheets and code. From this view the user moves on to questioning,
  visualising or mapping a library.}
  \label{fig:ui-libraries}
\end{figure}

\section{Grounded Chat}

The chat view, shown in Figure~\ref{fig:ui-chat}, is where a user asks questions
of a library. The answer is written in plain language and is followed by the
sources it relied on, so that each claim can be checked against the document it
came from. In the figure the user asks which student has the highest average, and
the platform answers with the specific value and lists the source it used. This is
the grounded, citing behaviour designed in Chapter~4 in action.

\begin{figure}[H]
  \centering
  \includegraphics[width=\textwidth]{QA_with_doc.png}
  \caption{The grounded chat. The user asks a question and the platform answers
  from the selected library, listing the source passages beneath the answer. The
  reply gives a specific, grounded result rather than a vague summary, and the
  cited source lets the user verify it. When the library cannot support an answer,
  the platform says so instead of guessing.}
  \label{fig:ui-chat}
\end{figure}

\section{Team Collaboration}

The same chat can be used by a team. Figure~\ref{fig:ui-team} shows the team view,
where members of an organisation share libraries and a conversation. Messages are
attributed to the member who sent them, and the assistant's grounded replies,
including tables where they help, are visible to everyone in real time. This turns
the platform from a personal tool into a shared workspace over a common set of
documents.

\begin{figure}[H]
  \centering
  \includegraphics[width=\textwidth]{Team_Chats.png}
  \caption{The team chat. Members of an organisation share libraries and a
  conversation, with each message attributed to its sender and the grounded
  replies visible to all in real time. Here a member asks for a comparison and the
  assistant responds with a structured answer. The realtime channel keeps every
  member's view in step.}
  \label{fig:ui-team}
\end{figure}

\section{Agent: Generating Visuals}

The agent view, shown in Figure~\ref{fig:ui-visuals}, turns a library or an
attached file into visuals. The user chooses the kinds of visual to allow, such as
bar, line or pie charts, and states a goal in plain language. The agent produces
the charts together with a short written explanation, including the assumptions it
made and recommendations for a clearer presentation, and each chart can be
downloaded as an image. The figure shows a grouped bar chart and a line chart
generated from a spreadsheet, with the numbers taken directly from the source.

\begin{figure}[H]
  \centering
  \includegraphics[width=\textwidth]{Generate_Visuals.png}
  \caption{The agent generating visuals. From a spreadsheet and a plain-language
  goal, the agent produces a grouped bar chart and a line chart, each downloadable
  as an image, alongside a written explanation with its assumptions and
  recommendations. Because the values are bound from the source file, the charts
  are faithful to the underlying data.}
  \label{fig:ui-visuals}
\end{figure}

\section{Agent: Generating Documents}

The agent also drafts written documents, as shown in Figure~\ref{fig:ui-docs}.
From a request and the connected sources it produces a document, offered both as a
readable draft and as a downloadable file in Markdown or PDF form. In the figure
the agent has drafted several documents from the user's computer-vision sources,
each available to read or download. This lets a user turn a library into a report
or a structured summary without leaving the platform.

\begin{figure}[H]
  \centering
  \includegraphics[width=\textwidth]{Generate_doc.png}
  \caption{The agent generating documents. From the connected sources the agent
  drafts written documents and offers them as downloadable Markdown or PDF files.
  Several drafts produced from the user's sources are shown, each ready to read or
  download. The documents are written from the library, so they reflect the user's
  own material.}
  \label{fig:ui-docs}
\end{figure}

\section{Saved Artefacts}

The artefacts the agent produces are kept for later use. Figure~\ref{fig:ui-saved}
shows the saved-visuals panel, which collects the charts a user has generated, each
with its title, a short description and a download option. This means that work
done with the agent is not lost when the conversation ends, and a user can return
to a chart and reuse it.

\begin{figure}[H]
  \centering
  \includegraphics[width=0.52\textwidth]{Save_Visual_docs.png}
  \caption{The saved-artefacts panel. Charts generated by the agent are collected
  here with their titles and descriptions, each downloadable as an image. Keeping
  the artefacts means a user can return to earlier work and reuse it rather than
  regenerate it. The panel gathers visuals produced across different requests in
  one place.}
  \label{fig:ui-saved}
\end{figure}

\section{Knowledge Graph}

Finally, the platform can map a library as a knowledge graph. Building the graph
runs as a background task with clear progress, shown in
Figure~\ref{fig:ui-kg-build}, as the system reads the documents and extracts
entities. The finished graph, shown in Figure~\ref{fig:ui-kg-view}, draws the
entities as nodes coloured by type, such as concept, method, system and product,
and lets the user explore how they connect across the library. The user can
rebuild the graph or switch to a different library.

\begin{figure}[H]
  \centering
  \includegraphics[width=\textwidth]{Knowledge_graph_generation.png}
  \caption{Building the knowledge graph. The platform reads the documents of a
  library and extracts entities and relationships as a background task, showing the
  number of passages processed and the overall progress. The user can cancel the
  build if needed. This is the construction flow from Chapter~4 as the user sees
  it.}
  \label{fig:ui-kg-build}
\end{figure}

\begin{figure}[H]
  \centering
  \includegraphics[width=\textwidth]{Knowledge_graph_representaiton.png}
  \caption{The interactive knowledge graph. The entities of a library are drawn as
  nodes coloured by type, such as concept, method, system and product, and the user
  can explore how they connect. A library selector and a rebuild control sit at the
  top. The graph turns a flat collection of documents into a structure the user can
  navigate.}
  \label{fig:ui-kg-view}
\end{figure}

\section{Summary}

This chapter presented the finished product through the web application. It showed
how a user connects and processes libraries, asks grounded questions and sees the
sources behind each answer, collaborates with a team in real time, generates
faithful visuals and written documents with the agent, keeps the resulting
artefacts, and maps a library as an interactive knowledge graph. These views
demonstrate that the platform delivers, in a single usable interface, the
capabilities set out at the start of the thesis. The next chapter evaluates how
well it performs.
